\ifmostrarGabarito
\begin{minted}{python}
from typing import List

def eh_primo(numero: int) -> bool:
    if numero <= 1:
        return False
    for divisor in range(2, int(numero**0.5) + 1):
        if numero % divisor == 0:
            return False
    return True

def filtrar_primos_intervalo(inicio: int, fim: int) -> List[int]:
    primos = []
    for num in range(inicio, fim + 1):
        if eh_primo(num):
            primos.append(num)
    return primos
\end{minted}

\textbf{Explicação:}
A solução adota a modularização em duas funções:
\begin{enumerate}
    \item \texttt{eh\_primo}: Testa se um número individual é primo, descartando números $\le 1$ e verificando divisores até $\sqrt{\text{numero}}$.
    \item \texttt{filtrar\_primos\_intervalo}: Percorre o intervalo fechado de \texttt{inicio} a \texttt{fim}, chamando \texttt{eh\_primo(num)} para filtrar e acumular os números primos em uma lista.
\end{enumerate}

\vspace{0.3cm}
\textbf{Critérios de Correção e Penalidades Sugeridas (Erros Comuns):}
\begin{itemize}[leftmargin=*, noitemsep, topsep=2pt]
    \item \textbf{Falta de modularização ou não chamar função auxiliar} (\textbf{--6,0 pontos}): Implementar tudo em uma única função ou não invocar \code{eh\_primo} dentro de \code{filtrar\_primos\_intervalo}.
    \item \textbf{Falha nos casos base $\le 1$} (\textbf{--5,0 pontos}): Não descartar números $\le 1$, classificando erroneamente $0$ ou $1$ como primos.
    \item \textbf{Intervalo de divisores incorreto} (\textbf{--5,0 pontos}): Omitir o $+1$ no teto da raiz quadrada (\code{range(2, int(n**0.5) + 1)}), falhando em quadrados perfeitos como 4, 9 e 25, ou iniciar em 1.
    \item \textbf{Retorno prematuro no laço} (\textbf{--5,0 pontos}): Colocar \code{return True} dentro de um bloco \code{else} na primeira iteração do laço divisor em \code{eh\_primo}.
    \item \textbf{Limites do intervalo ou retorno} (\textbf{--4,0 pontos}): Usar \code{range(inicio, fim)} esquecendo o limite superior inclusivo, ou não retornar a lista de primos resultante.
\end{itemize}
\else
\linhasResposta{18}
\clearpage
\linhasResposta[FOLHA DE CONTINUACAO / RASCUNHO -- QUESTAO Q9]{30}
\clearpage
\fi
